\documentclass[11pt]{article}
\usepackage[margin=1in]{geometry}
\usepackage{amsmath,amssymb,amsthm}
\usepackage{hyperref}
\usepackage{enumitem}
\usepackage{xurl}
\let\oldus\_
\renewcommand{\_}{\oldus\allowbreak}

\newtheorem{theorem}{Theorem}
\newtheorem{lemma}[theorem]{Lemma}
\newtheorem{corollary}[theorem]{Corollary}
\theoremstyle{definition}
\newtheorem{definition}[theorem]{Definition}
\theoremstyle{remark}
\newtheorem{remark}[theorem]{Remark}

\title{A disproof of Conjecture 00000004055\\[2pt]
\large Over the self-injective algebra $k[x]/(x^2)$ every syzygy period is $1$, so no prime period occurs}
\date{}

\begin{document}
\sloppy
\maketitle

\section{The conjecture}

Conjecture 00000004055 reads (English):
\begin{quote}
\emph{Definition:} The syzygy $\Omega(M)$ is the kernel of the minimal projective cover, and the
period is the minimal number of iterations first returning to an isomorphic module.
\emph{Conjecture:} On every self-injective algebra the syzygy period of any prescribed prime $p$
can be realized, and the minimal period 2 is attained by open-chain modules.
\end{quote}
The Chinese text says the same: on every self-injective algebra, a syzygy period equal to any
given prime $p$ can be realized.

\textbf{Reading.} Clause 1: for every self-injective algebra $R$ and every prime $p$ there is an
$R$-module $M$ of syzygy period $p$. Clause 2: every self-injective algebra has a module of
period $2$ (an ``open-chain module''; this notion is undefined, and we show that no module at
all has period $2$ in our example).

We refute both clauses with one algebra, $A=k[\varepsilon]$ with $\varepsilon^2=0$ (isomorphic to
$k[x]/(x^2)$), over any field $k$. It is a commutative $2$-dimensional algebra, and it is not
semisimple (the ideal $\varepsilon A$ is a nonzero nilpotent ideal).

\section{Definitions}

Let $R$ be a ring.
\begin{definition}
A submodule $K\le P$ is \emph{superfluous} if $K+L=P$ implies $L=P$ for every submodule $L$.
A \emph{projective cover} of $M$ is an epimorphism $\pi:P\to M$ with $P$ projective and $\ker\pi$
superfluous in $P$. ($R$ is \emph{self-injective} if $R$ is injective as an $R$-module.)
\end{definition}
These are the standard notions: Wikipedia, \emph{Projective cover}: ``A projective cover is a pair
(P,p), with P a projective object in'' $\mathcal{C}$ ``and p a superfluous epimorphism'', where for
$R$-modules ``a superfluous epimorphism is then an epimorphism'' $p:P\to M$ ``such that the kernel
of p is a superfluous submodule of P''. Wikipedia, \emph{Essential extension}: ``The dual notion of an
essential submodule is that of superfluous submodule (or small submodule).'' It defines them as in
Definition~1. Wikipedia, \emph{Injective
module}: ``If a ring is injective over itself as a right module, then it is called a right
self-injective ring.''

\begin{definition}[Lean: \texttt{IsSyzygy}, \texttt{IterSyz}, \texttt{HasPeriod}]
$N$ is a syzygy of $M$ if $N\cong\ker\pi$ for some projective cover $\pi:P\to M$.
$\Omega^0(M)\cong N$ means $M\cong N$, and $\Omega^{n+1}(M)\cong N$ means that there is a syzygy
$M'$ of $M$ with $\Omega^n(M')\cong N$. $M$ has \emph{period} $p$ if $p\ge1$,
$\Omega^p(M)\cong M$, and $\Omega^m(M)\not\cong M$ for $1\le m<p$.
\end{definition}
Projective covers are unique up to isomorphism when they exist, so this agrees with the usual
functional $\Omega$. We do not need that fact: Theorem~\ref{thm:main} shows that every syzygy of a
periodic module is isomorphic to it, so the result holds for any choice of covers. Modules are
arbitrary (not required to be finitely generated); this only weakens the existence claim.

\textbf{Conventions.} $\mathbb{N}$ contains $0$; periods are $\ge 1$. Clause 1 is formalized for
commutative algebras (\texttt{PrimePeriodsRealized}): for every commutative ring $R$ that is a
finite-dimensional $k$-algebra with $R$ injective over itself, and every prime $p$, some
$R$-module (in \texttt{Type}) has period $p$. For a commutative ring, left and right
self-injectivity coincide, and the clause for all self-injective algebras implies this one.

\section{Proof}

\begin{lemma}[\texttt{selfInjective}]
$A$ is self-injective.
\end{lemma}
\begin{proof}
Baer's criterion (Mathlib \texttt{Module.Baer.injective}): let $I$ be an ideal and $g:I\to A$
linear. If $I$ contains an element $x=a+b\varepsilon$ with $a\ne0$, then $x$ is a unit, $I=A$ and
$g$ extends to itself. Otherwise $I\subseteq\varepsilon A$. If $I=0$, take $0$. Otherwise
$x_0=b_0\varepsilon\in I$ with $b_0\ne0$; $c=g(x_0)$ satisfies $\varepsilon c=g(\varepsilon x_0)=0$, so
$c=c_1\varepsilon$, and $r\mapsto r\,(c_1/b_0)$ extends $g$.
\end{proof}

\begin{lemma}[\texttt{mem\_eps\_of\_eps\_smul}, \texttt{superfluous\_le\_eps}]
Let $P$ be a projective $A$-module. (a) If $\varepsilon p=0$ then $p\in\varepsilon P$. (b) Every
superfluous submodule $K$ of $P$ lies in $\varepsilon P$.
\end{lemma}
\begin{proof}
$P$ has a splitting $s:P\to A^{(P)}$ of the map $A^{(P)}\to P$ (Mathlib
\texttt{Module.projective\_def'}). If every coordinate of $s(p)$ lies in $\varepsilon A$, then
$s(p)=\varepsilon f$ and $p=\varepsilon\cdot(\text{image of }f)\in\varepsilon P$. (a): $\varepsilon s(p)=s(\varepsilon p)=0$
forces every coordinate into $\varepsilon A$. (b): if $y\in K\setminus\varepsilon P$, some coordinate
$\varphi(y)=s(y)_i$ is a unit. Let $L=\{p:\varphi(p)\in\varepsilon A\}$, a submodule with $y\notin L$.
Every $p$ equals $c y+(p-cy)$ with $c=\varphi(p)_0/\varphi(y)_0$ (constant terms), and $p-cy\in L$.
So $K+L=P$ but $L\ne P$, a contradiction.
\end{proof}

\begin{lemma}[\texttt{eps\_smul\_ker}, \texttt{ker\_equiv\_of\_eps}]
Let $\pi:P\to M$ be a projective cover. Then $\varepsilon$ kills $\ker\pi$. If $\varepsilon$ kills $M$, then
$\ker\pi\cong M$.
\end{lemma}
\begin{proof}
By (b), $\ker\pi\subseteq\varepsilon P$, and $\varepsilon^2=0$. If $\varepsilon M=0$, then $\mu:P\to\ker\pi$,
$p\mapsto\varepsilon p$, is well defined and onto (by (b)), and $\ker\mu=\ker\pi$: if $\varepsilon p=0$ then
$p=\varepsilon q$ by (a), so $\pi(p)=\varepsilon\pi(q)=0$; conversely $\ker\pi$ is killed by $\varepsilon$.
Hence $\ker\pi\cong P/\ker\mu=P/\ker\pi\cong M$.
\end{proof}

\begin{theorem}[\texttt{periodic\_imp\_period\_one}, \texttt{period\_eq\_one}, \texttt{no\_prime\_period}]\label{thm:main}
Let $M$ be an $A$-module with $\Omega^n(M)\cong M$ for some $n\ge1$. Then every syzygy of $M$ is
isomorphic to $M$, and $\Omega^1(M)\cong M$. Hence every module of finite period has period $1$, and
no $A$-module has period $p$ for any prime $p$.
\end{theorem}
\begin{proof}
Every syzygy is killed by $\varepsilon$ (\texttt{syzygy\_killed}), so $M\cong\Omega^n(M)$ is killed by
$\varepsilon$ (\texttt{iterSyz\_killed}). By the lemma, every syzygy of $M$ is isomorphic to $M$. If $M$
had period $p\ge2$, then $\Omega^1(M)\cong M$ would contradict minimality.
\end{proof}

\begin{theorem}[\texttt{simple\_has\_period\_one}, \texttt{minimal\_period\_is\_one}]
The simple module $\varepsilon A$ has period $1$: multiplication by $\varepsilon$, $A\to\varepsilon A$, is a
projective cover with kernel $\varepsilon A$. So the least period over $A$ is $1$, and no module has
period $2$.
\end{theorem}

\begin{corollary}[\texttt{conjecture\_4055\_false}, \texttt{period\_two\_false}]
For every field $k$, \texttt{PrimePeriodsRealized k} is false, and so is the period-2 clause
(every finite-dimensional commutative self-injective $k$-algebra has a module of period 2).
\end{corollary}

\section{Scope}

Refuted: the literal reading (for every self-injective algebra and every prime $p$ some module has
period $p$), already for $p=2$ and for every $p$; and the period-2 clause, whatever ``open-chain
module'' means. Not refuted: (a) the reading ``for every prime $p$ some self-injective algebra has a
module of period $p$''; (b) the reading in which $\Omega^n(M)\cong M$ is taken in the stable module
category (isomorphism up to projective summands). It is not formalized and we make no claim about
it, although Lemma~5 and \texttt{syzygy\_killed} show that every syzygy $N$ over $A$ already satisfies
$\Omega(N)\cong N$.

\section{Lean formalization}

File \texttt{lean/Conjecture4055/Basic.lean} (309 lines, Lean 4.33.1, Mathlib v4.33.1),
namespace \texttt{C4055}. Main theorem:
\texttt{C4055.conjecture\_4055\_false (k : Type) [Field k] : \textasciitilde PrimePeriodsRealized k}.
\texttt{\#print axioms} for \texttt{conjecture\_4055\_false}, \texttt{period\_two\_false},
\texttt{no\_prime\_period}, \texttt{periodic\_imp\_period\_one}, \texttt{minimal\_period\_is\_one}
and \texttt{selfInjective} reports only \texttt{propext}, \texttt{Classical.choice} and
\texttt{Quot.sound}. There is no \texttt{sorry} and no \texttt{native\_decide}.

\section*{Sources}
\begin{itemize}
\item \url{https://en.wikipedia.org/wiki/Projective_cover}
\item \url{https://en.wikipedia.org/wiki/Essential_extension}
\item \url{https://en.wikipedia.org/wiki/Injective_module}
\end{itemize}

\end{document}
